\section{结论}\label{sec:conclusion}

本文综述了费马大定理从1637年提出到1995年解决、再到其后续影响的研究历程。可以概括出三条主线：

\begin{enumerate}[label=(\arabic*)]
    \item \textbf{代数主线}：从欧拉、拉梅的逐指数硬攻，到库默尔的理想论——为了解决一个具体问题，数学家建立了一整套代数数论。FLT 是“好问题孕育好数学”的最佳例证。
    \item \textbf{几何主线}：弗雷曲线、塞尔 $\varepsilon$ 猜想与里贝特定理把丢番图方程翻译成椭圆曲线与伽罗瓦表示的语言；怀尔斯的模性提升与泰勒--怀尔斯配补完成了最后一跃。证明本身促进了朗兰兹纲领的形成与壮大。
    \item \textbf{技术遗产}：模性提升、欧拉系、水平约化如今是数论研究的基础设施，其后续成果（Serre 猜想、Sato--Tate 猜想、一般模性定理）的分量不亚于 FLT 本身。
\end{enumerate}

费马大定理的故事说明了数学的深刻统一性：一个关于整数的初等方程，其最终答案藏在模形式与伽罗瓦表示的对应之中。悬而未决的 Beal 猜想、广义费马方程与朗兰兹纲领的更一般对应，正是这一统一性的下一个战场。
